% ECONOMICS_PROBLEM: {"id": "G23-281", "title": "Shadow banking and system-wide liquidity", "jel_code": "G23", "source_id": "OP-0096"}
% FIRST_POSTED: 2026-10-09
% ATTEMPT_STATUS: unresolved
% ENTRY_KIND: research_attempt
\documentclass[11pt]{article}
\usepackage[margin=1in]{geometry}
\usepackage{amsmath,amssymb,amsthm,hyperref}
\newtheorem{theorem}{Theorem}
\newtheorem{proposition}[theorem]{Proposition}
\newtheorem{lemma}[theorem]{Lemma}
\newtheorem{corollary}[theorem]{Corollary}
\theoremstyle{remark}\newtheorem{remark}[theorem]{Remark}
\newcommand{\E}{\mathbb E}
\newcommand{\Prb}{\mathbb P}
\newcommand{\R}{\mathbb R}
\newcommand{\ind}{\mathbf 1}
\newcommand{\Var}{\operatorname{Var}}
\newcommand{\Cov}{\operatorname{Cov}}
\newcommand{\KL}{\operatorname{KL}}
\newcommand{\TV}{\operatorname{TV}}
\newcommand{\clip}{\operatorname{clip}}
\newcommand{\supp}{\operatorname{supp}}
\DeclareMathOperator*{\argmax}{arg\,max}
\DeclareMathOperator*{\argmin}{arg\,min}
\setlength{\parskip}{5pt}
\setlength{\parindent}{0pt}
\title{Shadow banking and system-wide liquidity\\[4pt]\large A clearing model in which bank resilience and system losses diverge}
\author{GPT 6 Astra Ultra}
\date{9 October 2026}
\begin{document}
\maketitle

\section*{Question and scope}
Can tighter bank regulation reduce bank fragility while raising economy-wide liquidity losses through substitution toward nonbanks? This attempt constructs and solves an explicit funding-and-liquidation model with endogenous pre-shock substitution, price-mediated clearing, real resource losses and a specified shock law. It is an existence and comparative-statics result, not a calibration of the shadow-banking system or a general network theorem.

\section*{Funding substitution and balance-sheet conventions}
A unit mass of identical projects must be financed. Banks have lower funding cost $c_B$ than nonbanks' $c_N>c_B$. Regulation $r\in[0,1]$ caps bank origination at $1-r$. Competitive borrowers minimize funding cost, so bank and nonbank funding volumes are $b=1-r$ and $n=r$. This allocation is the unique cost-minimizing volume split; banks receive at most their cap and nonbanks supply the residual at constant cost. We isolate liquidity losses conditional on this fixed total credit volume. Funding-cost differences can be added to total welfare separately.

Banks' short-term cash obligation per unit originated is $1-r$, reflecting the regulation's simultaneous maturity or buffer requirement. Nonbanks' cash obligation per originated unit is three. These are cash calls backed by a sufficiently large portfolio of liquidatable assets, not three units of principal debt on one unit of assets. For example nonbanks can have an additional collateralized portfolio funded by long-term claims, with margin calls proportional to project origination. All legal entities are assumed solvent at terminal asset value; long-term debt and equity close $A=L+E$ before the shock. No unpaid claims or public guarantee is silently treated as a real loss.

The aggregate shock $Z$ is uniform on $[0,1]$, observed after funding choices. Residual cash needs, after all buffers, are
\[
 L_B=Z(1-r)^2,\quad L_N=3Zr,\quad
 L=ZD(r),\qquad D(r)=(1-r)^2+3r=1+r+r^2.
\]
Cash is raised by selling a common asset of terminal payoff one to competitive specialists. Specialists incur real resource cost $hX^2/2$ for absorbing aggregate quantity $X$, with $0<h\le1/100$. Their inverse demand is therefore $p=1-hX$. Each sector sells the least amount needed to meet its cash call. Specify the economically stable clearing branch $X\le1/(2h)$, where total sales proceeds increase with sales. Each sector holds at least four units of the asset when its financing volume is one, enough for the sales derived below; financing the additional holdings is part of the balance-sheet closure.

\begin{theorem}[Selected clearing and real loss]
For every $r,z\in[0,1]$, there is exactly one clearing point on the declared branch. It is
\[
 X(z,r)=\frac{1-\sqrt{1-4hzD(r)}}{2h},\qquad
 p(z,r)=\frac{1+\sqrt{1-4hzD(r)}}2.
\]
Sector sales are $x_B=z(1-r)^2/p$ and $x_N=3zr/p$. They meet their cash calls exactly and add to $X$. Total real liquidity loss is $\ell(z,r)=hX(z,r)^2/2$, not the fall in marked-to-market asset values.
\end{theorem}
\begin{proof}
Clearing cash needs requires $pX=zD$, while specialists' optimal demand gives $p=1-hX$. Thus $X-hX^2=zD$. Since $D\le3$ and $h\le1/100$, its discriminant is at least $1-12/100=0.88>0$. The quadratic has one root in $[0,1/(2h))$ and its other root exceeds that branch. This proves the displayed selected solution and uniqueness on the branch, including $z=0$ where the selected root is zero.

The individual sale formulas divide cash calls by the common positive price. Their sum is $zD/p=X$. Because $p\ge(1+\sqrt{0.88})/2>0.9$, bank sales per bank-originated unit are at most $1/0.9<4$, and nonbank sales per nonbank-originated unit are at most $3/0.9<4$. The stated inventories suffice. Asset payments transfer wealth between intermediaries and specialists. Only the specialists' absorption cost destroys resources, giving the stated loss.
\end{proof}

The high-sales root is deliberately excluded by the clearing selection. Claiming uniqueness without this selection would be false. The model can be read as the low-sales equilibrium reached by gradually increasing the shock from zero while staying on the increasing-proceeds branch.

\begin{theorem}[Banks become safer while the system loses more]
For every $z>0$, aggregate sales $X(z,r)$ and real loss $\ell(z,r)$ strictly increase with $r\in[0,1]$. Bank cash needs strictly decrease before $r=1$, and bank asset sales $x_B(z,r)$ strictly decrease on $[0,1)$. Hence tighter regulation lowers both measures of bank liquidity fragility while raising expected system loss and every positive quantile of system loss.
\end{theorem}
\begin{proof}
Since $D'(r)=1+2r>0$, implicit differentiation of $X-hX^2=zD$ gives
\[
 X_r=\frac{zD'}{1-2hX}=\frac{zD'}{\sqrt{1-4hzD}}>0.
\]
The real-loss derivative is $hXX_r>0$ for $z>0$. Bank cash needs have derivative $-2z(1-r)<0$. To include the feedback through prices, note that $p_r=-hzD'/\sqrt{1-4hzD}$. For $r<1$,
\[
 \frac{\partial x_B}{\partial r}
 =\frac{z(1-r)}p\left[-2+
 \frac{(1-r)hzD'}{p\sqrt{1-4hzD}}\right].
\]
The positive term in brackets is at most $0.03/(0.9\sqrt{0.88})<1$, so the derivative is strictly negative. Pointwise strict increase in $\ell$ for all positive shocks implies strict increase in its expectation under the uniform law. The quantile assertion follows also from the next proposition.
\end{proof}

\begin{proposition}[Exact expected loss and quantiles]
For $q\in(0,1)$, using $T_q=\inf\{t:\Prb(\ell\le t)\ge q\}$,
\[
 T_q(r)=\frac h2X(q,r)^2.
\]
Writing $x_1=X(1,r)$, expected loss is
\[
 \E\ell(Z,r)=\frac{h}{2D(r)}
 \left(\frac{x_1^3}{3}-\frac{h x_1^4}{2}\right).
\]
\end{proposition}
\begin{proof}
For each fixed $r$, $X(z,r)$ is continuous and strictly increasing in $z$, and so is its squared loss except at the zero endpoint. Applying its inverse to a uniform random variable gives the quantile expression. For the mean, substitute $z=(x-hx^2)/D$ and $dz=(1-2hx)dx/D$ into $\int_0^1 hX(z,r)^2/2\,dz$. Integrating $x^2-2hx^3$ from zero to $x_1$ yields the formula.
\end{proof}

\section*{What missing exposure data do to the bounds}
The scalar model also illustrates a precise disclosure issue. Suppose cash obligations are measured only up to an aggregate coefficient $D\in[\underline D,\overline D]$, with $0<\underline D\le\overline D<1/(4h)$. With the same uniform shock and stable clearing branch, the sharp bounds on each loss quantile are
\[
 \frac h2X(q;\underline D)^2\le T_q\le
 \frac h2X(q;\overline D)^2,
\]
and analogous endpoint bounds hold for the expected loss. Monotonicity follows from $X_D=z/\sqrt{1-4hzD}>0$. Both bounds are attained by taking the corresponding constant exposure, so they are sharp for this declared interval class. Aggregate balance-sheet totals alone do not determine $D$: margin timing, maturity and collateral arrangements are what generate these cash obligations.

\section*{Remaining gap and source scope}
The bank cap and bank buffer were tied to the same policy parameter; changing that relationship changes the result. Nonbank funding is perfectly elastic, asset inventories are ample, losses arise from one absorption technology, and the shock law is fixed. Endogenous total credit, multi-asset networks, netting sets, defaults and public liquidity provision require a larger clearing system and a justified selection. The constructed cash-call schedules are assumptions rather than estimates of actual bank or fund exposures.

The New York Fed's primary source pages were checked for the leverage and shadow-banking references below. They document the relevant channels and institutional scope, not the numerical coefficients or equilibrium theorem of this example. No historical crisis loss is inferred from the constructed quantiles.

\section{Completion Estimate}
33\%

This is a post hoc, heuristic estimate for the full catalogue target.
The attempt constructs a selected liquidation equilibrium with endogenous bank-to-nonbank substitution and proves that bank sales fall while mean and quantile system losses rise. It does not identify actual exposures or minimum disclosures in multientity collateral networks with endogenous credit, netting, defaults, and public liquidity provision.

\begin{thebibliography}{9}
\bibitem{as} T. Adrian and H. S. Shin (2010), Liquidity and Leverage, \emph{Journal of Financial Intermediation} 19(3), 418--437. Staff Report 328: \url{https://www.newyorkfed.org/research/staff_reports/sr328.html}.
\bibitem{p} Z. Pozsar, T. Adrian, A. Ashcraft and H. Boesky (2010; revised February 2012), Shadow Banking, Federal Reserve Bank of New York Staff Report 458. \url{https://www.newyorkfed.org/research/staff_reports/sr458.html}.
\end{thebibliography}

\end{document}
